% GENERATED FILE. Edit canonical lessons, then run npm run generate:books.
% canonical-source-hash: fnv1a64:79661b09cb57bd59
% canonical-lessons: PT-C523-capitulo, PT-C523-romancista, PT-C523-biografia, PT-C523-reler, PT-C523-marcador

\chapter{Capítulo a capítulo}
\label{ch:pt-capitulo-a-capitulo}
\hlchaptermodality{523}

\begin{chapteropening}
\textbf{By the end of this chapter:} \emph{I can talk about novels, biographies and rereading.}

Complete the last lesson of chapter 523: I can talk about novels, biographies and rereading.
\end{chapteropening}

\section[o capítulo]{o capítulo — chapter}

\label{lesson:PT-C523-capitulo}



\begin{quote}
Before the new one: say the Portuguese for to sculpt, then the Portuguese for sculptor.
\end{quote}

\subsection*{You'll want to know: o capítulo}
\textbf{o capítulo} — \textquotedblleft{}chapter\textquotedblright{}.

\begin{quote}
Leio um capítulo antes de dormir.
\end{quote}

\begin{quote}
\emph{I read a chapter before I go to sleep.}
\end{quote}

Also an episode of a soap: \emph{o último capítulo}.

\begin{cousinweb}
\emph{capítulo} is a learned borrowing from Late Latin \emph{capitulum} \textquotedblleft{}a chapter\textquotedblright{} (EtymDB), a little \textquotedblleft{}head\textquotedblright{} of text, from \emph{caput} \textquotedblleft{}head\textquotedblright{}.
\end{cousinweb}

\begin{grammarlens}[title={What you've built}]
One more word for telling, explaining and describing at B1.
\end{grammarlens}

\subsection*{Your turn}
\begin{itemize}
\raggedright
  \item \emph{Say it:} \emph{o capítulo}
  \item \emph{Say it:} \emph{o capítulo}, once more
  \item \emph{Say it:} a sentence of your own with \emph{o capítulo}
  \item \emph{Recall:} say the Portuguese for to sculpt, then the Portuguese for sculptor, then say \emph{o capítulo} again
\end{itemize}

\begin{tcolorbox}[breakable,colback=teal!4,colframe=teal!35!black,title={Before you move on}]
What does o capítulo mean? (\textquotedblleft{}Chapter\textquotedblright{}.) Say it once more.
\end{tcolorbox}
\section[o romancista]{o romancista — novelist}

\label{lesson:PT-C523-romancista}



\begin{quote}
Before the new one: say the Portuguese for sculptor, then the Portuguese for chapter.
\end{quote}

\subsection*{You'll want to know: o romancista}
\textbf{o romancista} — \textquotedblleft{}novelist\textquotedblright{}.

\begin{quote}
O romancista apresentou o novo livro na feira.
\end{quote}

\begin{quote}
\emph{The novelist presented his new book at the fair.}
\end{quote}

Built on \emph{o romance} \textquotedblleft{}a novel\textquotedblright{}. A woman is \emph{a romancista}.

\begin{cousinweb}
\emph{romancista} is built on \emph{romance} \textquotedblleft{}a novel\textquotedblright{}, from Old Occitan \emph{romans} (EtymDB), a tale told in the everyday Romance tongue rather than in Latin.
\end{cousinweb}

\begin{grammarlens}[title={What you've built}]
One more word for telling, explaining and describing at B1.
\end{grammarlens}

\subsection*{Your turn}
\begin{itemize}
\raggedright
  \item \emph{Say it:} \emph{o romancista}
  \item \emph{Say it:} \emph{o romancista}, once more
  \item \emph{Say it:} a sentence of your own with \emph{o romancista}
  \item \emph{Recall:} say the Portuguese for sculptor, then the Portuguese for chapter, then say \emph{o romancista} again
\end{itemize}

\begin{tcolorbox}[breakable,colback=teal!4,colframe=teal!35!black,title={Before you move on}]
What does o romancista mean? (\textquotedblleft{}Novelist\textquotedblright{}.) Say it once more.
\end{tcolorbox}
\section[a biografia]{a biografia — biography}

\label{lesson:PT-C523-biografia}



\begin{quote}
Before the new one: say the Portuguese for chapter, then the Portuguese for novelist.
\end{quote}

\subsection*{You'll want to know: a biografia}
\textbf{a biografia} — \textquotedblleft{}biography\textquotedblright{}.

\begin{quote}
Estou a ler uma biografia de um pintor famoso.
\end{quote}

\begin{quote}
\emph{I am reading a biography of a famous painter.}
\end{quote}

\emph{Autobiografia} — a life written by its subject.

\begin{cousinweb}
\emph{biografia} is a learned word built from Greek \emph{bíos} \textquotedblleft{}life\textquotedblright{} and \emph{gráphein} \textquotedblleft{}to write\textquotedblright{}.
\end{cousinweb}

\begin{grammarlens}[title={What you've built}]
One more word for telling, explaining and describing at B1.
\end{grammarlens}

\subsection*{Your turn}
\begin{itemize}
\raggedright
  \item \emph{Say it:} \emph{a biografia}
  \item \emph{Say it:} \emph{a biografia}, once more
  \item \emph{Say it:} a sentence of your own with \emph{a biografia}
  \item \emph{Recall:} say the Portuguese for chapter, then the Portuguese for novelist, then say \emph{a biografia} again
\end{itemize}

\begin{tcolorbox}[breakable,colback=teal!4,colframe=teal!35!black,title={Before you move on}]
What does a biografia mean? (\textquotedblleft{}Biography\textquotedblright{}.) Say it once more.
\end{tcolorbox}
\section[reler]{reler — to reread}

\label{lesson:PT-C523-reler}



\begin{quote}
Before the new one: say the Portuguese for novelist, then the Portuguese for biography.
\end{quote}

\subsection*{You'll want to know: reler}
\textbf{reler} — \textquotedblleft{}to reread\textquotedblright{}.

\begin{quote}
Gosto de reler os livros da minha infância.
\end{quote}

\begin{quote}
\emph{I like rereading the books of my childhood.}
\end{quote}

Conjugated like \emph{ler}: \emph{releio}, \emph{relês}, \emph{relê}.

\begin{cousinweb}
\emph{reler} is built from \emph{re-} \textquotedblleft{}again\textquotedblright{} and \emph{ler} \textquotedblleft{}to read\textquotedblright{}, which is inherited from Latin \emph{legere} \textquotedblleft{}to gather, to read\textquotedblright{} (Lewis \& Short).
\end{cousinweb}

\begin{grammarlens}[title={What you've built}]
One more word for telling, explaining and describing at B1.
\end{grammarlens}

\subsection*{Your turn}
\begin{itemize}
\raggedright
  \item \emph{Say it:} \emph{reler}
  \item \emph{Say it:} \emph{reler}, once more
  \item \emph{Say it:} a sentence of your own with \emph{reler}
  \item \emph{Recall:} say the Portuguese for novelist, then the Portuguese for biography, then say \emph{reler} again
\end{itemize}

\begin{tcolorbox}[breakable,colback=teal!4,colframe=teal!35!black,title={Before you move on}]
What does reler mean? (\textquotedblleft{}To reread\textquotedblright{}.) Say it once more.
\end{tcolorbox}
\section[o marcador]{o marcador — bookmark; marker pen}

\label{lesson:PT-C523-marcador}



\begin{quote}
Before the new one: say the Portuguese for biography, then the Portuguese for to reread.
\end{quote}

\subsection*{You'll want to know: o marcador}
\textbf{o marcador} — \textquotedblleft{}bookmark; marker pen\textquotedblright{}.

\begin{quote}
Perdi o marcador e não sei em que página ia.
\end{quote}

\begin{quote}
\emph{I lost my bookmark and do not know what page I was on.}
\end{quote}

Also a felt-tip marker, and in football a goalscorer.

\begin{cousinweb}
\emph{marcador} is built on \emph{marcar} \textquotedblleft{}to mark\textquotedblright{}, which is usually traced to a Germanic word for a boundary sign.
\end{cousinweb}

\begin{grammarlens}[title={What you've built}]
One more word for telling, explaining and describing at B1.
\end{grammarlens}

\subsection*{Your turn}
\begin{itemize}
\raggedright
  \item \emph{Say it:} \emph{o marcador}
  \item \emph{Say it:} \emph{o marcador}, once more
  \item \emph{Say it:} a sentence of your own with \emph{o marcador}
  \item \emph{Recall:} say the Portuguese for biography, then the Portuguese for to reread, then say \emph{o marcador} again
\end{itemize}

\begin{tcolorbox}[breakable,colback=teal!4,colframe=teal!35!black,title={Before you move on}]
What does o marcador mean? (\textquotedblleft{}Bookmark; marker pen\textquotedblright{}.) Say it once more.
\end{tcolorbox}
